\documentclass[dvisvgm,tikz,border=3pt]{standalone}
\usepackage{amsmath,amssymb}
\usepackage{pgfplots}
\usepackage{CJKutf8}
\pgfplotsset{compat=1.18}
\usetikzlibrary{arrows.meta,calc,positioning}

\definecolor{themebg}{HTML}{30362d}
\definecolor{themefg}{HTML}{edf4e8}
\definecolor{themegray}{HTML}{9ea897}
\definecolor{themecurve}{HTML}{7eb6ff}
\definecolor{themealert}{HTML}{ff7b7b}
\definecolor{themeamber}{HTML}{f5b942}

\pgfdeclarelayer{background}
\pgfdeclarelayer{foreground}
\pgfsetlayers{background,main,foreground}

\begin{document}
\begin{CJK*}{UTF8}{gbsn}
\pagecolor{themebg}
\begin{tikzpicture}[color=themefg, text=themefg, >=Stealth]

% ── 阶段 1：初始几何状态 (双二次型共存) ──
\begin{scope}[shift={(0,0)}]
  \draw[themegray!70, line width=0.8pt, rounded corners=5pt] (-2.2,-2.3) rectangle (2.2,2.3);
  \node[font=\bfseries\small, text=themecurve, above] at (0, 2.4) {阶段1：初始几何状态};

  % 坐标轴
  \draw[->, themegray!50, line width=0.6pt] (-1.9, 0) -- (1.9, 0);
  \draw[->, themegray!50, line width=0.6pt] (0, -1.9) -- (0, 1.9);

  % 正定二次型 A (晶蓝倾斜椭球，倾角35度)
  \draw[themecurve, line width=2.0pt, rotate=35] (0,0) ellipse (1.6 and 0.8);
  \node[text=themecurve, font=\bfseries\scriptsize, right] at (0.8, 1.3) {椭球 $\boldsymbol{A} \succ 0$};

  % 一般对称矩阵 B (暖金虚线倾斜椭球，倾角 -25度)
  \draw[themeamber, dashed, line width=1.6pt, rotate=-25] (0,0) ellipse (1.8 and 0.6);
  \node[text=themeamber, font=\bfseries\scriptsize, below left] at (-0.3, -1.2) {二次型 $\boldsymbol{B}$};
\end{scope}

% ── 变换箭头 1：白化 P1 ──
\draw[->, themecurve, line width=2.4pt] (2.5, 0) -- 
  node[above=3pt, font=\bfseries\footnotesize, text=themecurve, fill=themebg, inner sep=1.5pt] {第一步：白化 $\boldsymbol{P}_1$} 
  node[below=3pt, font=\scriptsize, text=themegray, fill=themebg, inner sep=1.5pt] {$\boldsymbol{P}_1^T\boldsymbol{A}\boldsymbol{P}_1 = \boldsymbol{I}$} (4.3, 0);

% ── 阶段 2：白化后 (A 归一化为单位球，B 依然倾斜) ──
\begin{scope}[shift={(6.8,0)}]
  \draw[themegray!70, line width=0.8pt, rounded corners=5pt] (-2.2,-2.3) rectangle (2.2,2.3);
  \node[font=\bfseries\small, text=themecurve, above] at (0, 2.4) {阶段2：$\boldsymbol{A}$ 归一为单位球 $\boldsymbol{I}$};

  % 坐标轴
  \draw[->, themegray!50, line width=0.6pt] (-1.9, 0) -- (1.9, 0);
  \draw[->, themegray!50, line width=0.6pt] (0, -1.9) -- (0, 1.9);

  % 矩阵 A 变为正圆/单位球 (晶蓝)
  \draw[themecurve, line width=2.2pt] (0,0) circle (1.2);
  \node[text=themecurve, font=\bfseries\scriptsize, above] at (0, 1.3) {单位球 $\boldsymbol{I}$ (各向同性)};

  % 矩阵 B' = P1^T B P1 (暖金倾斜椭圆)
  \draw[themeamber, dashed, line width=1.6pt, rotate=-40] (0,0) ellipse (1.6 and 0.7);
  \node[text=themeamber, font=\bfseries\scriptsize, below right] at (0.3, -1.1) {$\boldsymbol{B}' = \boldsymbol{P}_1^T\boldsymbol{B}\boldsymbol{P}_1$};
\end{scope}

% ── 变换箭头 2：正交旋转 Q ──
\draw[->, themeamber, line width=2.4pt] (9.3, 0) -- 
  node[above=3pt, font=\bfseries\footnotesize, text=themeamber, fill=themebg, inner sep=1.5pt] {第二步：旋转 $\boldsymbol{Q}$} 
  node[below=3pt, font=\scriptsize, text=themegray, fill=themebg, inner sep=1.5pt] {$\boldsymbol{Q}^T\boldsymbol{Q} = \boldsymbol{I}$} (11.1, 0);

% ── 阶段 3：同时对角化完成 (A 保持单位球，B 旋转至主轴) ──
\begin{scope}[shift={(13.6,0)}]
  \draw[themegray!70, line width=0.8pt, rounded corners=5pt] (-2.2,-2.3) rectangle (2.2,2.3);
  \node[font=\bfseries\small, text=themeamber, above] at (0, 2.4) {阶段3：同时合同对角化完成};

  % 坐标轴
  \draw[->, themegray!50, line width=0.6pt] (-1.9, 0) -- (1.9, 0);
  \draw[->, themegray!50, line width=0.6pt] (0, -1.9) -- (0, 1.9);

  % 矩阵 A 依然是正圆/单位球 (旋转不破正圆)
  \draw[themecurve, line width=2.2pt] (0,0) circle (1.2);
  \node[text=themecurve, font=\bfseries\scriptsize, above] at (0, 1.3) {保持 $\boldsymbol{I}$ 不变};

  % 矩阵 B 完美落在直角坐标轴上 (标准形)
  \draw[themeamber, line width=2.2pt] (0,0) ellipse (1.6 and 0.7);
  \node[text=themeamber, font=\bfseries\scriptsize, below] at (0, -0.9) {对角阵 $\boldsymbol{\Lambda}$};
\end{scope}

% ── 底部代数构造链卡片 ──
\begin{scope}[shift={(0,-2.8)}]
  \node[draw=themegray!60, fill=themebg, rounded corners=4pt, line width=0.8pt, inner sep=8pt, text width=15.2cm, anchor=north west] {
    \centering
    {\bfseries\color{themecurve} 同时合同对角化构造代数链：$\boldsymbol{C} = \boldsymbol{P}_1\boldsymbol{Q}$} \\[4pt]
    \raggedright\small
    $\bullet$ \textbf{第一步（白化Whitening）}：因 $\boldsymbol{A} \succ 0$，取 $\boldsymbol{P}_1 = \boldsymbol{Q}_A\boldsymbol{\Lambda}_A^{-1/2} \implies \boldsymbol{P}_1^T\boldsymbol{A}\boldsymbol{P}_1 = \boldsymbol{I}$。\\
    $\bullet$ \textbf{第二步（正交主轴旋转）}：令实对称阵 $\boldsymbol{B}' = \boldsymbol{P}_1^T\boldsymbol{B}\boldsymbol{P}_1$，取正交阵 $\boldsymbol{Q}$ 使 $\boldsymbol{Q}^T\boldsymbol{B}'\boldsymbol{Q} = \boldsymbol{\Lambda}$。\\
    $\bullet$ \textbf{保持单位球不破坏}：由于 $\boldsymbol{Q}$ 正交，$\boldsymbol{Q}^T\boldsymbol{I}\boldsymbol{Q} = \boldsymbol{I}$，令 $\boldsymbol{C} = \boldsymbol{P}_1\boldsymbol{Q}$ 即有 $\boldsymbol{C}^T\boldsymbol{A}\boldsymbol{C} = \boldsymbol{I}$ 且 $\boldsymbol{C}^T\boldsymbol{B}\boldsymbol{C} = \boldsymbol{\Lambda}$。
  };
\end{scope}

\end{tikzpicture}
\end{CJK*}
\end{document}
